% BHU PYQ -- Auto-generated & error-corrected
\documentclass[11pt,a4paper]{article}

\usepackage[a4paper,top=2.5cm,bottom=2.5cm,left=2.8cm,right=2.8cm,headheight=14pt]{geometry}
\usepackage{amsmath,amssymb}
\usepackage[T1]{fontenc}
\usepackage[utf8]{inputenc}
\usepackage{lmodern}
\usepackage{microtype}
\usepackage{fancyhdr}
\usepackage{enumitem}
\usepackage{hyperref}
\usepackage{lastpage}
\usepackage{parskip}

\hypersetup{colorlinks=true,urlcolor=black,linkcolor=black,
  pdftitle={STA-201: Descriptive Statistics and Distribution Theory -- BHU PYQ},pdfauthor={Banaras Hindu University}}

\pagestyle{fancy}\fancyhf{}
\renewcommand{\headrulewidth}{0.4pt}\renewcommand{\footrulewidth}{0.4pt}
\fancyhead[L]{\small\textit{Banaras Hindu University}}
\fancyhead[R]{\small\textit{STA-201: Descriptive Statistics and Distribu}}
\fancyfoot[C]{\small Page \thepage\ of \pageref{LastPage}}

\newlist{parts}{enumerate}{2}
\setlist[parts,1]{label=(\alph*),leftmargin=*,itemsep=5pt,topsep=3pt}
\setlist[parts,2]{label=(\roman*),leftmargin=*,itemsep=3pt,topsep=2pt}
\newcommand{\pts}[1]{\hfill\textit{[\#1]}}

\begin{document}

\begin{center}
  {\large\bfseries BANARAS HINDU UNIVERSITY}\\[2pt]
  {\normalsize B.A. (Hons.) Semester II Examination, 2023-24}\\[6pt]
  \rule{0.8\linewidth}{0.5pt}\\[6pt]
  {\large\bfseries Paper No. STA-201: Descriptive Statistics and Distribution Theory}\\[4pt]
  {\normalsize Time: 3 Hours\hfill Maximum Marks: 70}
\end{center}
\rule{\linewidth}{0.4pt}
\smallskip
\noindent\textit{Instructions: Answer \textbf{five} questions including Question~1 which is compulsory. Symbols have their usual meanings.}
\bigskip

\noindent\textit{Instructions: Answer \textbf{five} questions including Question~1 which is compulsory. Symbols have their usual meanings.}
\bigskip

\medskip\noindent\textbf{Question 1.}
\begin{parts}
  \item Write a note on the use of a scatter diagram to obtain an idea of the extent and direction of the correlation coefficient. \pts{7}
  \item If $Y_1$, $Y_2$, and $Y_3$ are uncorrelated variables each having the same variance, then obtain the correlation coefficient between $(Y_1+Y_2)$ and $(Y_2+Y_3)$. \pts{7}
\end{parts}
\medskip\rule{\linewidth}{0.2pt}

\medskip\noindent\textbf{Question 2.}
\begin{parts}
  \item Define intra-class correlation and also derive the expression for the intra-class correlation coefficient. \pts{7}
  \item Describe the concepts of multiple and partial correlation coefficients. Show that the multiple correlation coefficient $R_{1.23}$ in the usual notations is given by $R_{1.23}^2 = 1 - \frac{\Delta}{\Delta_{11}}$. \pts{7}
\end{parts}
\medskip\rule{\linewidth}{0.2pt}

\medskip\noindent\textbf{Question 3.}
\begin{parts}
  \item Define the consistency of data for attributes and derive the necessary conditions of consistency up to three attributes. \pts{7}
  \item Find the coefficient of correlation between the ages of the following mothers ($X$) and daughters ($Y$):
  \begin{center}
    \begin{tabular}{r|cccccc}
      Mother's Age ($X$ in years) & 24 & 26 & 28 & 30 & 32 & 34 \\ \hline
      Daughter's Age ($Y$ in years) & 4 & 5 & 6 & 8 & 9 & 11
    \end{tabular}
  \end{center} \pts{7}
\end{parts}
\medskip\rule{\linewidth}{0.2pt}

\medskip\noindent\textbf{Question 4.}
\begin{parts}
  \item Prove that $|E(X)| \le E(|X|)$, provided all expectations exist. \pts{7}
  \item Calculate the mathematical expectation of the sum of points when rolling a fair die. \pts{7}
\end{parts}
\medskip\rule{\linewidth}{0.2pt}

\medskip\noindent\textbf{Question 5.}
\begin{parts}
  \item Show that $\text{Var}(X) = E[\text{Var}(X|Y)] + \text{Var}[E(X|Y)]$, where all notations have their usual meaning. \pts{7}
  \item Define the moment generating function and state its important properties. \pts{7}
\end{parts}
\medskip\rule{\linewidth}{0.2pt}

\medskip\noindent\textbf{Question 6.}
\begin{parts}
  \item Define the geometric distribution and prove that it satisfies the lack of memory property. If $\bar{X}$ is the sample mean of $n$ i.i.d. random variables, show that $M_{\bar{X}}(t) = [M_X(t/n)]^n$. \pts{7}
  \item If $X_i \sim \text{Gamma}(a, b_i)$ for $i=1, 2, \dots, n$ are independent random variables, then prove that $\sum_{i=1}^n X_i \sim \text{Gamma}(a, \sum_{i=1}^n b_i)$. \pts{7}
\end{parts}
\medskip\rule{\linewidth}{0.2pt}

\medskip\noindent\textbf{Question 7.}
\begin{parts}
  \item If $X$ is a random variable and $E[X^2] < \infty$, then prove that $P(|X| \ge a) \le \frac{E[X^2]}{a^2}$ for all $a > 0$. \pts{7}
  \item Let the joint probability density function of random variables $X$ and $Y$ be:
  \[
  f(x, y) = e^{-(x+y)}, \quad x > 0, \ y > 0
  \]
  Find the following:
  \begin{parts}
    \item $P(X > 1)$
    \item $P(1 < X+Y < 2)$
    \item $P(X < Y \mid X < 2Y)$
  \end{parts} \pts{7}
\end{parts}
\medskip\rule{\linewidth}{0.2pt}

\medskip\noindent\textbf{Question 8.}
Describe any \textbf{three} of the following in brief: \pts{14}
\begin{parts}
  \item Rank Correlation
  \item Normal Distribution
  \item Intra-Class Correlation
  \item Uniform Distribution
\end{parts}

\vfill
\rule{\linewidth}{0.4pt}
\begin{center}\small
  \href{https://bhu-exam.vercel.app/}{Click here to visit our website}\\[4pt]\href{https://me-aryan.vercel.app/}{Click here to visit our contributor}\\[4pt]\href{https://quashan-me.vercel.app/}{Click here to visit our contributor}
\end{center}

\end{document}